\section{Hardware}\label{sec:hardware}

\begin{table}[t]
\centering\small
\caption{The test machine. Everything in this report was measured on this one host.}\label{tab:host}
\begin{tabularx}{\linewidth}{@{}lX@{}}
\toprule
GPUs & 2 $\times$ NVIDIA GeForce RTX 5060 Ti 16\,GB (Blackwell, compute capability 12.0), 180\,W power
limit each, 16{,}311\,MiB usable per card, \textbf{no NVLink, no peer-to-peer path} \\
GPU attachment & GPU0 on a CPU PCIe root port; GPU1 behind the chipset bridge, sharing its uplink with
the SATA controller (both disks and the swap file), the network interfaces and USB. Both links train
at x4 on x8-capable cards \\
CPU / RAM & AMD Ryzen 5 8500G, 6 cores / 12 threads; \textbf{14\,GiB} usable system RAM; 4\,GiB swap
(12\,GiB during the second battery, \cref{sec:corrections}) \\
Storage & Two SATA disks: system (218\,GB) and a separate 110\,GB model store \\
Driver & NVIDIA 595.84, CUDA 13.2; Ubuntu Linux; all engines run in Docker containers \\
Power sensing & GPU power from the driver at 1\,Hz; CPU package from RAPL; \emph{no} wall-socket, BMC
or IPMI sensor. ``System'' power anywhere in this report is modelled, not measured \\
\bottomrule
\end{tabularx}
\end{table}

\begin{figure}[t]
\centering
\begin{tikzpicture}[font=\footnotesize,
  box/.style={draw=black!70,rounded corners=2pt,minimum height=0.85cm,minimum width=2.5cm,align=center,fill=white},
  gpu/.style={box,fill=cbsky!25,minimum width=3.3cm,minimum height=1.25cm},
  >=Stealth]
\node[box,fill=cborange!25,minimum width=3.0cm,minimum height=1.1cm] (cpu) {CPU (6C/12T)\\14\,GiB RAM};
\node[gpu,right=2.6cm of cpu,yshift=1.1cm] (g0) {\textbf{GPU0} 16\,GB\\layers $0..k$ + their KV};
\node[box,below=1.35cm of cpu] (pch) {chipset bridge};
\node[gpu,right=2.6cm of pch] (g1) {\textbf{GPU1} 16\,GB\\layers $k{+}1..63$ + KV\\output head, draft context};
\node[box,left=1.0cm of pch,minimum width=2.1cm] (io) {SATA (disks, swap)\\NIC, USB};
\draw[<->,line width=1pt] (cpu.east) -- node[above,sloped]{PCIe x4 (root port)} (g0.west);
\draw[<->,line width=1pt] (cpu.south) -- node[right]{shared uplink} (pch.north);
\draw[<->,line width=1pt] (pch.east) -- node[above]{PCIe x4} (g1.west);
\draw[<->] (pch.west) -- (io.east);
\draw[<->,dashed,cbred,line width=0.9pt] (g0.south) -- node[right,text=cbred,align=left,font=\scriptsize,pos=0.45]{no P2P: activations\\cross via host} (g1.north);
\end{tikzpicture}
\caption{The two cards are not a memory pool and are not attached symmetrically. With layer split,
each layer's weights \emph{and its slice of the KV cache} live on one card, so a configuration fails
when the fuller card runs out, however much is free on the other. GPU1 also carries the output head
and the speculative draft context, and sits behind the shared chipset uplink.}\label{fig:topology}
\end{figure}

\Cref{tab:host} lists the machine; \cref{fig:topology} shows the part of it that matters most. Three
properties shape everything that follows.

\textbf{The cards are two 16\,GB budgets, not one 32\,GB budget.} There is no NVLink and the driver
reports no peer-to-peer path. Of the split modes llama.cpp offers, only layer split
(\flag{-sm layer}) runs on the engine builds that survive (\cref{sec:servers}); it places whole layers,
with their KV slices, on one card or the other. The binding limit is therefore 16{,}311\,MiB on
whichever card is fuller. \Cref{sec:context} shows default placement stranding 2--3.8\,GiB on GPU0 while
GPU1 ran out.

\textbf{The cards are attached asymmetrically.} GPU1 shares a chipset uplink with both disks, the swap
file and the network. I did not test whether this contributes to the per-card asymmetries measured
later; the documented cause of those is layer placement. It is recorded because a host description
that omits it would be incomplete.

\textbf{System RAM is small, and it binds twice.} With the model memory-mapped, about 10\,GiB is
available. This (i)~caps the standard KL-divergence tool at 8{,}192 tokens of context, because it holds
a chunk's logits in memory at 248\,KB per token (at 16{,}384 tokens the host fell to 305\,MiB free and
wrote nothing in 13 minutes), and (ii)~caused the study's only serving outage
(\cref{sec:hostram}). No full-precision reference fits anywhere: BF16 weights are 54.7\,GB, so every
divergence in this report is measured against the largest quantized build, not against the original
model.

\section{Software versions and model builds}\label{sec:software}

\begin{table}[t]
\centering\small
\caption{Engines and tools. A result is tied to the image that produced it; image digests are in the
repository's \texttt{env-manifest.json}.}\label{tab:software}
\begin{tabularx}{\linewidth}{@{}l>{\raggedright\arraybackslash}p{4.4cm}X@{}}
\toprule
label & version & used for \\
\midrule
\textbf{E-A} llama.cpp (MTP build) & 0.3.0-dev, commit \texttt{d222767}; image \texttt{feb0231976b6} &
first battery: divergence, HellaSwag, HumanEval+, RULER, split sweeps, MTP speculation \\
\textbf{E-B} llama.cpp DFlash fork & 0.1.2-dev build 50, commit \texttt{f7aadef}; image \texttt{22bb8b7fed8b} &
the only DFlash2-capable build during the first battery \\
\textbf{E-C} llama.cpp upstream & build b10975, commit \texttt{4c9233c}; image \texttt{5268283a8d65} &
second battery: both drafters on one build; KV map, speed, AA-LCR, GPQA, soak; now serving \\
\textbf{E-H} llama.cpp PR\,\#27342 build & commit \texttt{1deefcc} (\emph{image deleted}) &
all measurements before 2026-08-29, marked \hist{} \\
vLLM & 0.27.1 (stable) and a 0.26.1rc1 nightly & NVFP4 accuracy; NVFP4 speculative speed \\
SGLang & 0.5.18 & attempted; never started (\cref{sec:servers}) \\
EvalPlus & patched local install & HumanEval / HumanEval+ scoring \\
mini-SWE-agent & 2.4.6 & SWE-bench Verified agent scaffold \\
RULER & reference generators and metric & long-context retrieval \\
\bottomrule
\end{tabularx}
\end{table}

\begin{table}[t]
\centering\small
\caption{Model builds. All are Unsloth Dynamic GGUFs of Qwen3.8-27B and all embed the MTP head.
\QsixXL{} is the divergence reference. Sizes in bytes as downloaded; SHA-256 prefixes are in the
repository.}\label{tab:arms}
\shrinktofit{\begin{tabular}{@{}llrrl@{}}
\toprule
build & role & bytes & GiB & batteries \\
\midrule
\QsixXL & divergence reference & 25{,}299{,}061{,}664 & 23.56 & first \\
\Qsix   & arm & 21{,}983{,}677{,}344 & 20.47 & first, second \\
\Qfive  & arm & 20{,}876{,}938{,}144 & 19.44 & first \\
\Qfour  & arm & 17{,}559{,}178{,}144 & 16.35 & first, second \\
\IQfour & historical only & 14{,}252{,}845{,}984 & 13.27 & \hist \\
\Qthree, \arm{Q4\_K\_M} & historical only (files deleted) & --- & --- & \hist \\
NVFP4 (W4A4, FP8 KV) & vLLM checkpoint & $\approx$22\,GB & --- & \hist \\
DFlash2 drafter \arm{Q4\_K\_M} & speculative sidecar & 1{,}143{,}006{,}816 & 1.06 & second \\
\bottomrule
\end{tabular}}
\end{table}

\Cref{tab:software,tab:arms} pin the stack. Two facts about versions turned out to matter more than
expected.

\textbf{Deleting one image froze the exploratory period's results.} Build E-H was the only one on which tensor
split (\flag{-sm tensor}) ever worked on these cards. On E-A and E-B the flag is listed in
\flag{--help} and the server dies with a CUDA illegal-memory-access error. When E-H was deleted during
routine disk cleanup, every tensor-split result and every 262K-token result measured until then became
unreproducible, and two model files from that period were also gone. Re-measuring on surviving builds
gave \emph{different} context ceilings. Those early rows are kept, marked \hist{}, and never placed
in a table with later ones. \emph{Flag availability is not capability.}

\textbf{A drafter file is bound to an engine build.} The DFlash2 drafter loads only on builds that
know its tensor layout. On E-A it fails with \texttt{wrong number of tensors; expected 81, got 58},
and a harness that scored by exit status recorded a clean 0.000 pass@1, indistinguishable in a results
table from a model that ran and failed. The same signature later blocked a production cut-over
(\cref{sec:config}). The drafter file itself was also revised upstream during the study (a rope-metadata
fix); the first-battery DFlash2 rows used the earlier revision on a different engine from the MTP
rows, which is why \cref{sec:spec} ranks the drafters only from the second battery.

\paragraph{Engine defaults are not what the documentation says.} On E-A, a request with no sampling
fields runs at temperature 1.0, top-$k$ 20, top-$p$ 0.95, min-$p$ 0.05. That matches neither
llama.cpp's documented defaults (0.8 / 40 / 0.95 / 0.05) nor either of Qwen's presets. The server
echoes effective settings per request, so this is cheap to check; every run here sends an explicit
sampling block and asserts the echo. The model also \emph{reasons by default}. The widely quoted
\texttt{/no\_think} prompt suffix does not disable it on this chat template. What does is the request field
\begin{center}\verb|chat_template_kwargs: {"enable_thinking": false}|\end{center}
while \verb|{"reasoning_effort": "none"}| raises a template exception (accepted values are
\texttt{low}, \texttt{medium}, \texttt{xhigh}). A two-token probe distinguishes the states in about a second and
ran before every battery.
